\documentclass[12pt, a4paper]{article}
\usepackage[T2A]{fontenc}
\usepackage[utf8]{inputenc}
\usepackage[russian]{babel}
\usepackage{amsmath, amssymb} % Для математических символов
\usepackage{enumitem} % Для гибкой настройки списков
\usepackage{geometry} % Для настройки полей
\geometry{top=2cm, bottom=2cm, left=2.5cm, right=2cm}

\begin{document}

\section*{Георгий Игоревич Вольфсон <<Делимость с человеческим лицом>>}

\subsection*{Задача 19 на профильном ЕГЭ по математике}

\subsubsection*{Задача 1}
\textbf{На доске написано несколько различных натуральных чисел, которые делятся на 3 и оканчиваются на 4.}
\textbf{а) Может ли их сумма составлять 282?}
\textbf{б) Может ли их сумма составлять 390?}
\textbf{в) Какое наибольшее количество чисел могло быть на доске, если их сумма равна 2226?}

\textbf{Решение:}
Числа делятся на 3 и оканчиваются на 4. Пусть число имеет вид \(10m + 4\). Тогда \(10m + 4 \equiv m + 1 \equiv 0 \pmod{3}\), откуда \(m \equiv 2 \pmod{3}\). Значит, числа имеют вид \(10(3k+2) + 4 = 30k + 24\), где \(k \ge 0\). Последовательность: \(24, 54, 84, 114, \ldots\)

\textbf{а)} Да, например: \(24 + 54 + 204 = 282\). Все числа делятся на 3, оканчиваются на 4 и различны.

\textbf{б)} Нет. Сумма \(n\) чисел равна \(30K + 24n = 390\). Тогда \(30K = 390 - 24n\). При \(n=5\) сумма пяти наименьших чисел \(24+54+84+114+144 = 420 > 390\), значит \(n \le 4\). Проверяем:
\begin{itemize}
    \item \(n=1\): \(390 - 24 = 366\) — не делится на 30.
    \item \(n=2\): \(390 - 48 = 342\) — не делится.
    \item \(n=3\): \(390 - 72 = 318\) — не делится.
    \item \(n=4\): \(390 - 96 = 294\) — не делится.
\end{itemize}
Сумма 390 невозможна.

\textbf{в)} Сумма \(n\) чисел: \(30K + 24n = 2226\), откуда \(30K = 2226 - 24n\). Правая часть должна делиться на 30. \(2226 \equiv 6 \pmod{30}\), \(24n \equiv 24n \pmod{30}\). Нужно \(6 - 24n \equiv 0 \pmod{30} \Rightarrow 24n \equiv 6 \pmod{30} \Rightarrow 4n \equiv 1 \pmod{5} \Rightarrow n \equiv 4 \pmod{5}\).

Сумма \(n\) наименьших чисел: \(15n^2 + 9n\). При \(n=12\): \(2268 > 2226\). При \(n=11\): \(1914 \le 2226\). Значит, \(n \le 11\). Возможные \(n \equiv 4 \pmod{5}\) при \(n \le 11\): \(n=4\) и \(n=9\). Максимальное \(n=9\).

Проверим \(n=9\): \(30K = 2226 - 216 = 2010 \Rightarrow K = 67\). Возьмём \(k_i = 0,1,2,3,4,5,6,7,39\) (сумма 67). Числа: \(24, 54, 84, 114, 144, 174, 204, 234, 1194\). Сумма 2226.

\textbf{Ответ:} а) да; б) нет; в) 9.

\subsubsection*{Задача 2}
\textbf{На доске было написано несколько различных натуральных чисел. Эти числа разбили на три группы, в каждой из которых оказалось хотя бы одно число. К каждому числу из первой группы приписали справа цифру 6, к каждому числу из второй группы приписали справа цифру 9, а числа третьей группы оставили без изменений.}
\textbf{а) Могла ли сумма всех этих чисел увеличиться в 9 раз?}
\textbf{б) Могла ли сумма всех этих чисел увеличиться в 19 раз?}
\textbf{в) В какое наибольшее количество раз могла увеличиться сумма всех этих чисел?}

\textbf{Решение:}
Пусть в первой группе \(x\) чисел с суммой \(S_x\), во второй — \(y\) чисел с суммой \(S_y\), в третьей — \(z\) чисел с суммой \(S_z\). Исходная сумма \(S = S_x + S_y + S_z\). Новая сумма:
\[
S' = (10S_x + 6x) + (10S_y + 9y) + S_z.
\]

\textbf{а)} Да. Пример: группы \((1)\), \((2, 5)\), \((4)\). Исходная сумма \(1+2+5+4=12\). Новая: \(16 + 29 + 59 + 4 = 108 = 9 \cdot 12\).

\textbf{б)} Нет. Если \(S' = 19S\), то:
\[
10S_x + 6x + 10S_y + 9y + S_z = 19(S_x + S_y + S_z).
\]
\[
2x + 3y = 3S_x + 3S_y + 6S_z.
\]
Так как \(S_x \ge x\) и \(S_y \ge y\), то \(2x+3y \ge 3x+3y \Rightarrow x \le 0\), что противоречит \(x \ge 1\).

\textbf{в)} Оценим \(\frac{S'}{S}\). Так как \(T \ge c\) (в третьей группе \(c\) натуральных чисел), \(a \ge 1\), \(c \ge 1\):
\[
6a + 9b - 9T \le 6a + 9b - 9c = 9(a+b+c) - 3a - 18c \le 9n - 21.
\]
Кроме того, так как все \(n\) чисел различны, \(S \ge \frac{n(n+1)}{2}\). Следовательно:
\[
\frac{S'}{S} \le 10 + \frac{2(9n - 21)}{n(n+1)}.
\]
Рассмотрим неравенство \(\frac{2(9n-21)}{n(n+1)} \le \frac{8}{5}\). Оно равносильно \(4n^2 - 41n + 105 \ge 0\), то есть \((n-5)(4n-21) \ge 0\). Это верно при \(n \le 5\) и \(n \ge 6\) (так как \(21/4 = 5.25\), для целых \(n \ge 6\) выполнено).
*Проверим \(n=5\):* \(\frac{2(45-21)}{30} = \frac{48}{30} = 1.6 = \frac{8}{5}\). Для \(n \ge 6\) правая часть убывает, так что оценка сохраняется.
Значит, \(\frac{S'}{S} \le 10 + \frac{8}{5} = \frac{58}{5} = 11.6\).

Достижимость: возьмём числа \(1, 2, 3, 4, 5\) (сумма 15). Третья группа — \(\{1\}\) (сумма \(T=1\)); вторая — \(\{2, 3, 4\}\) (получаем 29, 39, 49); первая — \(\{5\}\) (получаем 56). Новая сумма:
\[
1 + 29 + 39 + 49 + 56 = 174.
\]
Отношение \(\frac{174}{15} = 11.6 = \frac{58}{5}\).

\subsubsection*{Задача 3}
\textbf{На доске написано несколько различных натуральных чисел, в записи которых могут быть только цифры 1 и 6.}
\textbf{а) Может ли сумма этих чисел быть равна 173?}
\textbf{б) Может ли сумма этих чисел быть равна 109?}
\textbf{в) Какое наименьшее количество чисел может быть на доске, если их сумма равна 1021?}

\textbf{Решение:}
Заметим, что любое число, состоящее из цифр 1 и 6, оканчивается на 1 или 6, а значит, даёт остаток 1 при делении на 5.

\textbf{а)} Да, например: \(166 + 6 + 1 = 173\).

\textbf{б)} Нет. Сумма \(n\) таких чисел даёт остаток \(n \pmod{5}\). Число 109 даёт остаток 4 при делении на 5, значит \(n \equiv 4 \pmod{5}\), откуда \(n = 4, 9, 14, \ldots\)
При \(n \ge 9\) сумма девяти наименьших чисел \(1+6+11+16+61+66+111+116+161 = 549 > 109\), что невозможно.
Остаётся \(n=4\). Если среди чисел есть 66, то сумма остальных трёх равна \(109-66=43\), но максимальная сумма трёх различных чисел из допустимых, не превосходящих 43 (это 1, 6, 11, 16), равна \(16+11+6=33 < 43\). Если есть 61, то остальные три в сумме дают 48, что также больше 33. Если нет ни 61, ни 66, то сумма четырёх чисел не превосходит \(16+11+6+1=34 < 109\). Противоречие.

\textbf{в)} Сумма 1021 даёт остаток 1 при делении на 5, значит \(n \equiv 1 \pmod{5}\), откуда \(n = 1, 6, 11, \ldots\)
Одно число не подходит, так как в записи 1021 есть цифры 0 и 2.
Покажем, что \(n=6\) достижимо. Пример:
\[
661 + 161 + 166 + 16 + 11 + 6 = 1021.
\]
Все числа различны и состоят только из цифр 1 и 6.

\textbf{Ответ:} а) да; б) нет; в) 6.

\subsubsection*{Задача 4}
\textbf{Склад представляет собой прямоугольный параллелепипед с целыми сторонами, контейнеры — прямоугольные параллелепипеды с размерами $1 \times 1 \times 3$ м. Контейнеры на складе можно класть как угодно, но параллельно границам склада.}
\begin{enumerate}[label=\alph*)]
    \item Может ли оказаться, что полностью заполнить склад размером 120 кубометров нельзя?
    \item Может ли оказаться, что на склад объемом 100 кубометров не удастся поместить 33 контейнера?
    \item Пусть объем склада равен 800 кубометров. Какой процент объема такого склада удастся гарантированно заполнить контейнерами при любой конфигурации склада?
\end{enumerate}

\textbf{Решение.}

\paragraph{a)} Такого случая быть не может: склад объёмом $120\text{ м}^3$
всегда можно заполнить полностью.

Пусть размеры склада равны $a\times b\times c$, где $a,b,c$ --- целые
числа. Тогда
\[
abc=120.
\]
Поскольку $120$ делится на $3$, хотя бы одно из чисел $a,b,c$ делится на $3$.
Пусть, например, $c=3k$. Разобьём склад на $ab$ параллелепипедов размером
$1\times1\times c$. Каждый из них заполняется $k$ контейнерами,
расположенными вдоль стороны длины $c$. Следовательно, весь склад можно
заполнить полностью.

\textbf{Ответ:} нет, не может.

\paragraph{b)} Рассмотрим склад размером
\[
1\times2\times50.
\]
Его объём равен $100\text{ м}^3$. Контейнер можно расположить только вдоль
стороны длины $50$, поскольку две другие стороны короче $3$ м.

Таким образом, склад разбивается на две независимые части размером
$1\times1\times50$. В каждой помещается не более
\[
\left\lfloor\frac{50}{3}\right\rfloor=16
\]
контейнеров. Значит, всего на складе можно разместить не более $32$
контейнеров, а $33$ контейнера разместить нельзя.

\textbf{Ответ:} да, например, для склада $1\times2\times50$.

\paragraph{c)} Пусть размеры склада равны $a\times b\times c$, причём
\[
abc=800.
\]
Запишем каждую сторону в виде
\[
a=3q_1+r_1,\qquad b=3q_2+r_2,\qquad c=3q_3+r_3,
\]
где $r_1,r_2,r_3\in\{0,1,2\}$.

Разобьём каждую сторону на отрезки длины $3$ и, возможно, один остаток
длины $r_i$. Эти разбиения делят весь склад на прямоугольные
параллелепипеды. Каждый такой параллелепипед, у которого хотя бы одна
сторона имеет длину $3$, полностью заполняется контейнерами,
направленными вдоль этой стороны. Поэтому незаполненным может остаться
только параллелепипед размером
\[
r_1\times r_2\times r_3.
\]

Поскольку $800$ не делится на $3$, ни одна из сторон $a,b,c$ не делится
на $3$, то есть все $r_i$ ненулевые. Кроме того,
\[
r_1r_2r_3\equiv abc\equiv800\equiv2\pmod3.
\]
При $r_1,r_2,r_3\in\{1,2\}$ возможны лишь два случая:
\[
(r_1,r_2,r_3)=(1,1,2)
\quad\text{с точностью до перестановки},
\]
либо
\[
r_1=r_2=r_3=2.
\]
В первом случае останутся незаполненными $2\text{ м}^3$, во втором ---
$8\text{ м}^3$. Следовательно, всегда можно заполнить не менее
\[
800-8=792\text{ м}^3,
\]
что составляет
\[
\frac{792}{800}\cdot100\%=99\%
\]
объёма склада.

Покажем, что больший процент гарантировать нельзя. Рассмотрим склад
размером $2\times2\times200$. Все контейнеры должны быть направлены вдоль
стороны длины $200$. Склад разбивается на четыре части размером
$1\times1\times200$, в каждой из которых помещается не более
\[
\left\lfloor\frac{200}{3}\right\rfloor=66
\]
контейнеров. Поэтому всего можно разместить не более
\[
4\cdot66=264
\]
контейнеров. Их общий объём равен
\[
264\cdot3=792\text{ м}^3.
\]
Такое размещение действительно существует, а незаполненным остаётся
куб размером $2\times2\times2$.

Таким образом, гарантированно можно заполнить ровно
\[
\boxed{99\%}
\]
объёма склада.

\subsubsection*{Задача 5}
\textbf{Есть синие и красные карточки. Всего карточек 50 штук. На каждой карточке написано натуральное число. Среднее арифметическое всех чисел равно 16. Все числа на синих карточках разные. При этом любое число на синей карточке больше, чем любое число на красной. Числа на синих карточках увеличили в 2 раза, после чего среднее арифметическое стало равно 31,2.}
\begin{enumerate}[label=\alph*)]
    \item Может ли быть 10 синих карточек?
    \item Может ли быть 10 красных карточек?
    \item Какое наибольшее количество синих карточек может быть?
\end{enumerate}

\textbf{Решение:}
Пусть $S_{red}$ — сумма чисел на красных карточках, $S_{blue}$ — сумма чисел на синих.
Изначально сумма всех чисел:
\[
S = 16 \cdot 50 = 800.
\]
После увеличения синих чисел вдвое среднее стало 31,2:
\[
S' = 31,2 \cdot 50 = 1560.
\]
Сумма увеличилась только за счёт синих чисел (каждое увеличилось на само себя), значит прирост равен $S_{blue}$:
\[
S_{blue} = 1560 - 800 = 760.
\]
Тогда сумма чисел на красных карточках:
\[
S_{red} = 800 - 760 = 40.
\]

\textbf{а)} Если синих карточек 10, то красных 40. Сумма красных чисел равна 40. Так как все они натуральные, каждое из них равно 1. 
Тогда любое синее число должно быть больше 1 (то есть $\ge 2$). Синих 10 штук, все числа разные, их сумма 760. 
Например, можно взять числа: $2, 3, 4, 5, 6, 7, 8, 9, 10, 706$. Их сумма равна 760.
Условия задачи выполнены.
\textbf{Ответ:} да, может.

\textbf{б)} Если красных карточек 10, то синих 40. 
Сумма красных чисел равна 40. По принципу Дирихле среди 10 натуральных чисел, дающих в сумме 40, хотя бы одно число не меньше 4. Значит, максимальное красное число не меньше 4.
Так как любое синее число строго больше любого красного, минимальное синее число не меньше 5.
Тогда сумма 40 различных синих чисел будет не меньше суммы первых 40 натуральных чисел, начиная с 5:
\[
S_{\min} = 5 + 6 + \ldots + 44 = \frac{(5 + 44) \cdot 40}{2} = 980.
\]
Но по условию сумма всех синих чисел равна 760. Так как $980 > 760$, такое невозможно.
\textbf{Ответ:} нет, не может.

\textbf{в)} Пусть синих карточек $k$, тогда красных $50-k$. 
Сумма красных чисел равна 40. Значит, максимальное из чисел на красных карточках не меньше $\lceil 40/(50-k) \rceil$ (по принципу Дирихле). Следовательно, минимальное синее число не меньше $M + 1$, где $M = \lceil 40/(50-k) \rceil$.

Чтобы сумма $k$ различных синих чисел была как можно меньше (и не превышала 760), они должны идти подряд, начиная с этого минимального значения. Сумма первых $k$ таких чисел равна:
\[
S_{\min} = \frac{k(2(M+1) + k - 1)}{2} = \frac{k(2M + k + 1)}{2}.
\]
Эта сумма не должна превосходить 760.

Проверим возможные значения $k$:
\begin{itemize}
    \item Если $k = 36$ (тогда красных $50-36=14$): $M = \lceil 40/14 \rceil = 3$. 
    \[
    S_{\min} = \frac{36(2 \cdot 3 + 36 + 1)}{2} = \frac{36 \cdot 43}{2} = 774 > 760.
    \]
    Значит, 36 синих карточек быть не может.
    
    \item Если $k = 35$ (тогда красных $50-35=15$): $M = \lceil 40/15 \rceil = 3$. 
    \[
    S_{\min} = \frac{35(2 \cdot 3 + 35 + 1)}{2} = \frac{35 \cdot 42}{2} = 735 \le 760.
    \]
    Это возможно: например, возьмём 10 красных карточек с числом 1 и 5 красных с числом 2 (их сумма 20). Максимальное красное число — 2. Тогда синие числа начинаются с 3. Сумма 35 синих чисел от 3 до 37 равна 700. Оставшиеся 60 можно добавить к самому большому синему числу. Условия задачи выполнены.
\end{itemize}

\textbf{Ответ:} 35.


\end{document}